\chapter{What a language model actually does}
\label{ch:what-lm}

\newthought{A language model} is a function that takes a sequence of symbols
and returns a probability distribution over what symbol could come next. That
is the entire idea. Everything in this book — embeddings, attention, scaling
laws, KV caches, subquadratic alternatives — is engineering on top of that
one definition.

\section{The one equation you need}

Let \(x_1, x_2, \ldots, x_n\) be a sequence of \term{tokens} (Chapter~\ref{ch:tokens}
defines that word precisely; for now, think ``approximately a word''). A
language model is a parameterised distribution
\begin{equation}
  p_\theta(x_1, x_2, \ldots, x_n)
  \;=\;
  \prod_{t=1}^{n} p_\theta\!\left(x_t \,\middle|\, x_1, \ldots, x_{t-1}\right).
  \label{eq:chain-rule}
\end{equation}
The right-hand side is the chain rule of probability.\sidenote{This is the
\emph{autoregressive} factorisation. Predict each token given the prefix; multiply
the per-token probabilities to get the joint. Almost every modern LLM is
trained this way.} The model's job is to learn the conditional
\(p_\theta(x_t \mid x_{<t})\) — the probability of the \emph{next} token given
everything that came before.

Training is just maximum-likelihood estimation: pick parameters \(\theta\) that
make a large corpus of real text look as likely as possible.\sidenote{Equivalently,
minimise the negative log-likelihood, which (averaged over tokens) is what
people report as ``cross-entropy loss''. Lower is better. \(\log\) is natural
unless someone says otherwise.}

\keyidea{An LM is a next-token probability machine. Training is fitting that
machine to a corpus by maximum likelihood. Almost everything else is plumbing
that makes the machine fast and accurate.}

\section{Why that humble idea is so powerful}

A model that can predict ``the cat sat on the \rule{1cm}{0.4pt}'' has, in
order to do so, picked up a startling amount of structure. To put a high
probability on \emph{mat} (or \emph{floor}, or \emph{lap}) and a low
probability on \emph{microservice}, the model has to encode:

\begin{itemize}
\item the syntax of English noun phrases,
\item the typical co-occurrence of cats and household objects,
\item the rough fact that ``sat on'' takes a flat surface as its object,
\item and probably a sense of register (this is a children's-book sentence,
not a stack-trace).
\end{itemize}

None of that was hand-coded. It is a side effect of the maximum-likelihood
objective on enough text. \emph{Generation} is then the same model run
forwards: sample a token from the distribution, append it, and repeat. That
loop is the entire reason ChatGPT works.\sidenote{The other half is
instruction tuning and reinforcement learning from human feedback, which we
will not cover. The pretraining objective in
equation~\eqref{eq:chain-rule} is what builds the latent capability; RLHF
shapes the surface behaviour.}

\section{What the input \emph{actually} looks like}

Equation~\eqref{eq:chain-rule} talks about tokens as if they were given. They
are not — there is no canonical way to chop text into tokens, and the choice
matters. We unpack that in Chapter~\ref{ch:tokens}. For now: imagine each
\(x_t\) as one element of a finite \emph{vocabulary} of size~\(V\), so the
model's output at each step is a vector of \(V\) probabilities that sum
to~one.\sidenote{Typical \(V\) is between \(30{,}000\) and \(200{,}000\) for a
modern LLM. GPT-2 used 50,257; Llama 3 uses 128,000. Larger vocabularies
shorten the average input but make the output projection more expensive.}

\section{What the output \emph{actually} looks like}

The model produces a vector \(z_t \in \R^{V}\) of \term{logits} — pre-softmax
scores — and turns them into a probability distribution by
\begin{equation}
  p_\theta(x_t = v \mid x_{<t})
  \;=\;
  \frac{\exp(z_{t,v} / \tau)}{\sum_{v'=1}^{V} \exp(z_{t,v'} / \tau)},
  \label{eq:softmax}
\end{equation}
where \(\tau > 0\) is the \term{temperature}. At training time \(\tau = 1\)
and the loss is the negative log probability of the true next token. At
generation time you can sample directly, take the argmax (greedy decoding),
keep only the top-\(k\) most probable tokens (top-\(k\) sampling), or restrict
to the smallest set whose probability exceeds \(p\) (\term{nucleus} or
top-\(p\) sampling).\sidenote{Lower \(\tau\) sharpens the distribution and
makes the model deterministic-looking. Higher \(\tau\) makes it surprising,
sometimes unhelpfully so. Most product systems use \(\tau\) between
\(0.2\) and \(1.0\) and combine it with top-\(p\)~\(\approx\)~\(0.9\).}

\section{Why we care about \emph{how} \(p_\theta\) is computed}

The chain-rule view in equation~\eqref{eq:chain-rule} is architecture-free. It
does not say whether \(p_\theta(x_t \mid x_{<t})\) should be computed by a
hash table, an n-gram count, an LSTM, or a transformer. Those are
engineering choices, and they have very different cost profiles in compute
and in memory. The rest of this book is about that choice — how the field
arrived at the transformer, how that choice scales, and what is being
proposed in 2026 to replace it.

In particular, the cost of computing \(p_\theta(x_t \mid x_{<t})\) is not a
fixed constant per token. It depends on \emph{how much context} \(x_{<t}\) the
model is asked to attend to. For an n-gram model the cost is constant
(Chapter~\ref{ch:pre-transformer}). For an RNN it is also roughly constant
per token, but information from far back is lost (also
Chapter~\ref{ch:pre-transformer}). For a dense transformer it grows
\emph{quadratically} with context length (Chapter~\ref{ch:why-quadratic}).
That last fact is what the second half of the book is about.

\fillme{FILL-01-01}{Worked example: cross-entropy on a toy vocabulary}{%
  Construct a 4-token vocabulary \(\{a, b, c, d\}\) and a 5-token training
  sequence (e.g.\ \texttt{a b a c d}). Show by hand the negative
  log-likelihood under (i) the uniform distribution, (ii) the empirical
  unigram distribution from the sequence, and (iii) a perfect bigram model
  fit to the sequence. Discuss in two sentences why (iii) trivially wins on
  in-distribution data and why this overfits. End with a note on what
  per-token cross-entropy means in nats vs bits.}{%
  Cover and Thomas, \emph{Elements of Information Theory}, ch.~5;
  \texttt{research/notes.md \S 1}.}{%
  Half a page including the table of probabilities; one short equation per
  case.}

\fillme{FILL-01-02}{Code listing: a 30-line next-token sampler}{%
  Provide a self-contained Python listing (no PyTorch, no transformers
  library) that takes a function
  \texttt{logits\_fn(prefix: list[int]) -> list[float]} as an oracle and
  implements: (a) greedy decoding, (b) temperature sampling, (c) top-\(k\),
  (d) nucleus / top-\(p\). Use the \texttt{lstlisting} environment with
  \texttt{language=Python}. Comment each function with one line on \emph{when}
  to use it. Don't import \texttt{numpy} — keep it pure Python so the listing
  reads as runnable pseudocode.}{%
  Equation~\ref{eq:softmax}; the Holtzman et al.\ nucleus-sampling paper if
  citing top-\(p\).}{%
  ${\sim}30$ lines of code, plus 2 sentences of intro and 1 sentence of
  outro.}
